\section{An execution contract for grounded retrieval}\label{sec:C16:contract}
A proposed grounded-retrieval contract has six stages: define the information need, enforce access, retrieve candidates, verify source/version identity, assemble evidence with roles, and validate the resulting claim or action. Each stage emits inspectable artifacts. The contract is not a universal algorithm and does not guarantee that a model reasons correctly. It specifies the interfaces at which errors and unsupported transitions can be detected.

For execution, augment the evidence packet with freshness requirements and a precondition check. A retrieved resource state may have changed between search and tool dispatch. If the action depends on that state, revalidate it or use a conditional write supported by the remote service. Retrieval confidence cannot substitute for concurrency control. A perfect passage about an old state can justify the wrong present action.

Index compression and evidence compression should be audited separately. A fixed-dimensional document vector can discard the distinction needed to select a source, while a later passage packer can discard that distinction after the correct source was selected. Increasing the generation window repairs neither failure automatically. In a constructed diagnostic, first supply the correct document identifier while preserving the ordinary passage selector; then supply the decisive passage while preserving the ordinary packer. The two interventions locate losses at different boundaries. They are oracle diagnostic conditions, not deployable baselines. Account for the document-search, passage-scoring, and generation costs independently, and retain the final evidence spans so that the apparent benefit of a larger index cannot be confused with a larger model-visible evidence budget.

For scientific synthesis, retain the distinction among author-reported results, independently reproduced results, deductions, and proposed experiments. A retrieved repository can establish that code exists at a commit; it does not establish that its published experiments reproduce. An official product document can establish a documented interface; it does not reveal every proprietary implementation decision. The output should carry the same evidence boundary that the source supports.

Current 2026 retrieval mechanisms demonstrate several directions: query-conditioned runtime construction, associative traversal, shared procedural continuations, and calibrated influence of retrieved atoms. They solve different parts of the chain. JAM's preserved raw pages address pre-query information loss, MRAgent's active access addresses search adaptation, MATM's state keys address procedural matching, and MemCalib addresses use after evidence is supplied. Combining them is a design hypothesis, not an established composite result.

The original synthesis is to view retrieval as a controlled transition from stored evidence to active decision state. The transition must preserve source identity, role, time, and access scope as well as content. Its success is measured by justified behavior under bounded cost, not by the number of records returned. Chapter~\ref{ch:C17} places this transition within a unified consistency model linking the stored history, assembled context, policy, tools, and environment.
